\chapter{Analysis and Insights}
\label{ch:analysis}

This chapter examines four cross-cutting analytical questions: how vulnerability knowledge
transfers across languages at the CWE level, what kinds of errors the model makes on
real-world code, how well-calibrated its probability outputs are, and which CWE classes
show the largest performance gap between synthetic and real-world settings.

\section{Insight 1 --- Cross-Language Transfer at the CWE Level}

The aggregated F1 matrix in Experiment~2 showed that C-only models transfer to Python
better than Python-only models transfer to C.
Decomposing the ablation results by individual CWE class reveals why this asymmetry exists
and which concepts genuinely transfer versus which are language-locked.

\begin{table}[H]
  \centering
  \caption{Shared CWE F1 under cross-language evaluation (selected classes).}
  \label{tab:cwe_transfer}
  \begin{tabular}{lcccccc}
    \toprule
    CWE & C-only$\to$C & C-only$\to$Py & Py-only$\to$C & Py-only$\to$Py & Joint$\to$C & Joint$\to$Py \\
    \midrule
    CWE-798 & 1.00 & 1.00 & 1.00 & 1.00 & 1.00 & 1.00 \\
    CWE-22  & 1.00 & 1.00 & 0.50 & 1.00 & 1.00 & 1.00 \\
    CWE-327 & 1.00 & 0.90 & 1.00 & 0.93 & 1.00 & 1.00 \\
    CWE-89  & 0.50 & 0.92 & 0.90 & 1.00 & 0.63 & 1.00 \\
    CWE-134 & 0.50 & 0.90 & 0.50 & 1.00 & 0.53 & 1.00 \\
    CWE-732 & 0.40 & 1.00 & 0.38 & 1.00 & 0.67 & 1.00 \\
    CWE-78  & 1.00 & 1.00 & 0.17 & 1.00 & 0.70 & 1.00 \\
    CWE-367 & 1.00 & 0.00 & 0.47 & 0.00 & 1.00 & 0.00 \\
    \bottomrule
  \end{tabular}
\end{table}

\begin{figure}[H]
  \centering
  \begin{subfigure}[b]{0.48\textwidth}
    \includegraphics[width=\textwidth]{insight1_cwe_transfer_c.png}
    \caption{CWE transfer F1 on C test set.}
  \end{subfigure}
  \hfill
  \begin{subfigure}[b]{0.48\textwidth}
    \includegraphics[width=\textwidth]{insight1_cwe_transfer_python.png}
    \caption{CWE transfer F1 on Python test set.}
  \end{subfigure}
  \caption{Per-CWE F1 heatmaps for all three model variants across both test languages.}
  \label{fig:insight1}
\end{figure}

Three distinct transfer profiles emerge.

\textbf{Bidirectional transfer (CWE-798, CWE-22).}
Hardcoded credentials and path traversal transfer perfectly in both directions under all
three model variants.
These CWEs have structurally identical surface forms across C and Python: a string literal
assigned directly to a password variable, or a raw string path constructed by concatenating
user input.
The model learns the same token-level pattern regardless of the language it was trained on.

\textbf{Asymmetric transfer (CWE-78, CWE-89, CWE-134).}
Command injection (CWE-78) illustrates the strongest asymmetry.
The C-only model achieves F1\,=\,1.00 on Python command injection because Python's
\texttt{os.system()} and \texttt{subprocess.call()} share token patterns with C's
\texttt{system()}.
Conversely, the Python-only model scores only F1\,=\,0.17 on C command injection:
Python training provides no exposure to the C pattern involving buffer preparation with
\texttt{sprintf} followed by \texttt{system()}.
A similar asymmetry appears for CWE-89: the Python-only model transfers to C at 0.90
because string concatenation patterns in database queries look similar in both languages,
but the C-only model reaches only 0.50 on C SQL injection, reflecting that C database
code (raw \texttt{sqlite3}/\texttt{mysql\_query} APIs) differs substantially from
Python ORM-style queries.

\textbf{Language-locked (CWE-367, CWE-732).}
TOCTOU race conditions (CWE-367) appear at F1\,=\,0.00 for all models on Python test
examples: the Python realisation of a TOCTOU (\texttt{os.path.exists()} before file open)
is an underrepresented idiom relative to the C \texttt{stat()} / \texttt{open()} sequence.
CWE-732 shows the opposite distribution: both C-only and Python-only models score high on
Python permission errors (F1\,=\,1.00) because Python's \texttt{os.chmod(path, 0o777)}
is a compact, distinctive operation, but both models score poorly on C (F1\,$\approx$\,0.38--0.40)
because C file permission management spreads across multiple API calls without a single
identifying token.

The key insight is that transfer quality correlates with the degree to which a CWE's surface
representation is shared across languages --- not with the semantic similarity of the
underlying flaw.

\section{Insight 2 --- Real-World Error Analysis}

The real-world confusion matrix (TP\,=\,88, FP\,=\,31, TN\,=\,53, FN\,=\,28) provides
two distinct failure signals: false negatives (missed vulnerabilities) and false positives
(false alarms).

\subsection{False Negatives}

The 28 false negatives are concentrated in a small number of CWE classes.
CWE-787 (out-of-bounds write) and CWE-416 (use-after-free) each account for four misses;
CWE-400 (uncontrolled resource consumption), CWE-703 (improper exception handling),
CWE-284 (improper access control), and CWE-264 (permissions/access control) each contribute
three.
The majority of these CWEs are absent from the synthetic training set entirely.
CWE-787 does not appear in training at all; CWE-703 involves missing exception handling
rather than a dangerous function call, making it invisible to a model trained on patterns
of dangerous API usage.

\begin{figure}[H]
  \centering
  \includegraphics[width=0.75\textwidth]{insight2_missed_cwes.png}
  \caption{Distribution of false negatives (missed vulnerabilities) by CWE class on the
           real-world test set.}
  \label{fig:insight2_missed}
\end{figure}

\subsection{False Positives}

The 31 false positives cluster around CWE-703 (6 false alarms), CWE-787 (5), and CWE-125 (4).
These are CWEs where real-world safe code shares structural features with dangerous code:
error handling routines that check return values from memory operations look superficially
like code that performs unsafe memory operations.

\subsection{Function Length vs.\ Accuracy}

\begin{table}[H]
  \centering
  \caption{Real-world accuracy by function length.}
  \label{tab:length_accuracy}
  \begin{tabular}{lcc}
    \toprule
    Length bucket & Lines & Accuracy \\
    \midrule
    Short  & 1--30 & 0.709 \\
    Medium & 31--80 & 0.581 \\
    Long   & 81+   & 0.846 \\
    \bottomrule
  \end{tabular}
\end{table}

\begin{figure}[H]
  \centering
  \includegraphics[width=0.6\textwidth]{insight2_length_vs_accuracy.png}
  \caption{Real-world classification accuracy as a function of function length.}
  \label{fig:insight2_length}
\end{figure}

Accuracy is lowest on medium-length functions (31--80 lines).
Short real-world functions are often utility wrappers with a clear verdict.
Long functions tend to contain a single highly distinctive dangerous call that is easy
to identify even within a long body.
Medium-length functions are the hardest: long enough to contain confounders, but short enough
that any single dangerous-looking line carries disproportionate weight in the sequence
representation.

\section{Insight 3 --- Calibration}

A well-calibrated model is one whose predicted probability accurately reflects actual
likelihood.
Label smoothing with $\varepsilon = 0.1$ is the primary calibration mechanism: by replacing
hard targets with soft targets, the training objective penalises the model for pushing
probabilities to the extremes, producing moderate-confidence outputs for uncertain examples.

Following training, temperature scaling is applied.
The optimal temperature minimising ECE on the synthetic validation set is $T = 0.981$
--- so close to 1.0 that it produces negligible change in model outputs.
This directly confirms that label smoothing has already done most of the calibration work:
models trained without label smoothing typically require temperatures of 1.2--1.8 to achieve
comparable calibration.

\begin{figure}[H]
  \centering
  \begin{subfigure}[b]{0.48\textwidth}
    \includegraphics[width=\textwidth]{insight3_calibration_synthetic.png}
    \caption{Calibration on synthetic held-out set.}
  \end{subfigure}
  \hfill
  \begin{subfigure}[b]{0.48\textwidth}
    \includegraphics[width=\textwidth]{insight3_calibration_real_world.png}
    \caption{Calibration on real-world test set.}
  \end{subfigure}
  \caption{Reliability diagrams comparing predicted probability to empirical accuracy.}
  \label{fig:calibration}
\end{figure}

The synthetic calibration plot shows a characteristic bimodal output distribution:
probabilities clustered near 0 (clearly safe) and near 0.9--1.0 (clearly vulnerable),
with very few examples in the 0.2--0.8 uncertain range.
The real-world calibration plot shows a broader distribution with more probability mass
in the intermediate range, indicating that real functions present more ambiguous features ---
partial matches to vulnerability patterns that the model has not seen in combination before.

\section{Insight 4 --- Per-CWE Generalisation Gap}

Table~\ref{tab:gap_cwe} shows the comparison between held-out synthetic recall and
real-world recall for the ten CWEs that appear in both evaluation sets with sufficient
representation.

\begin{table}[H]
  \centering
  \caption{Per-CWE recall gap (synthetic held-out vs.\ real-world).}
  \label{tab:gap_cwe}
  \begin{tabular}{lllcc}
    \toprule
    CWE & Description & Held-out Recall & Real-world Recall & Gap \\
    \midrule
    CWE-732 & Insecure file permissions & 0.750 & 1.000 & $+0.250$ \\
    CWE-78  & Command injection         & 0.880 & 1.000 & $+0.120$ \\
    CWE-89  & SQL injection             & 0.878 & 1.000 & $+0.122$ \\
    CWE-120 & Buffer overflow           & 1.000 & 1.000 & $\phantom{+}0.000$ \\
    CWE-190 & Integer overflow          & 1.000 & 1.000 & $\phantom{+}0.000$ \\
    CWE-193 & Off-by-one error          & 1.000 & 1.000 & $\phantom{+}0.000$ \\
    CWE-22  & Path traversal            & 1.000 & 1.000 & $\phantom{+}0.000$ \\
    CWE-476 & Null pointer dereference  & 1.000 & 0.800 & $-0.200$ \\
    CWE-416 & Use-after-free            & 0.667 & 0.429 & $-0.238$ \\
    CWE-134 & Format string inj.\       & 0.767 & 0.000 & $-0.767$ \\
    \bottomrule
  \end{tabular}
\end{table}

\begin{figure}[H]
  \centering
  \includegraphics[width=0.75\textwidth]{insight4_gap_per_cwe.png}
  \caption{Per-CWE recall gap between the synthetic held-out set and the real-world test set.
           Positive values indicate better real-world performance.}
  \label{fig:insight4}
\end{figure}

The results fall into three groups.

\textbf{Positive transfer (CWE-732, CWE-78, CWE-89).}
Three CWEs perform better in the real world than on the synthetic held-out set.
All three are expressed through a small number of highly canonical API calls ---
\texttt{system()}, \texttt{execl()}, SQL string concatenation with \texttt{sprintf} ---
that appear consistently across real-world projects.
Real-world examples may be more archetypal than some held-out synthetic examples,
which occasionally embed the pattern within contextual wrappers that obscure the dangerous
operation.

\textbf{Zero gap (CWE-120, CWE-190, CWE-193, CWE-22).}
Four CWEs show perfect recall in both settings.
These are classes where the vulnerable operation is so syntactically distinctive ---
\texttt{strcpy} into a fixed array, unchecked arithmetic, path concatenation without
normalisation --- that real-world examples present essentially the same token pattern
as the training examples.

\textbf{Negative transfer (CWE-476, CWE-416, CWE-134).}
CWE-476 drops by 0.20: real-world null dereferences arise from complex call chains where
the pointer originates from an external API return value.
CWE-416 worsens to 0.429 from an already-low 0.667, consistent with its status as the
hardest class in the synthetic evaluation.
CWE-134 collapses to zero in the real world: the synthetic training contains Python template
injection examples, but DiverseVul's real-world CWE-134 instances are C format string
vulnerabilities --- \texttt{printf(user\_string)} accepting a user-controlled format
argument --- a C-specific idiom structurally different from the Python template injection
patterns in training.

The overall pattern across all four insights is coherent: CWEs expressed through compact,
canonical, language-consistent API calls generalise well; CWEs requiring inter-statement
reasoning or expressed differently across code corpora show the largest gaps.

\section{Insight 5 --- Token Attribution Analysis}
\label{sec:insight5}

The semantic understanding tests in Section~\ref{sec:understanding} established
\textit{whether} predictions are correct; token attribution analysis addresses
\textit{why} by quantifying the contribution of individual tokens to the model's output.

\subsection{Method}

Occlusion-based attribution replaces each non-special token in the input sequence with
the tokeniser's \texttt{[PAD]} token and records the resulting change in predicted
probability.
The attribution score for token $t_i$ is defined as:
\[
  s_i = p_{\text{baseline}} - p_{\text{ablated}(i)}
\]
where $p_{\text{baseline}}$ is the model's predicted vulnerability probability on the
original sequence and $p_{\text{ablated}(i)}$ is the probability after replacing $t_i$
with \texttt{[PAD]}.
A positive score indicates the token pushes the prediction toward VULNERABLE; a negative
score indicates it pushes toward SAFE.
The method is model-agnostic and makes no assumptions about the internal attention
structure; a high score proves the model relies on that token, but a near-zero score
does not prove the token is ignored --- it may indicate that the prediction is
distributed across many tokens simultaneously.

\subsection{Results by CWE Class}

Three qualitatively distinct attribution patterns emerge across the test cases.

\textbf{Correct attribution (CWE-416).}
For the use-after-free example, \texttt{free} receives the highest attribution score
(0.506), \texttt{req} --- the pointer that is subsequently dereferenced --- scores
second (0.060), and the \texttt{->} dereference operator also registers (0.003).
The model has assigned the largest weights to the semantically correct tokens: the
deallocation call and the identifier that connects the freed memory to its illegal
subsequent access.
This is the strongest evidence in the attribution analysis that the model has learned a
structural pattern rather than an incidental surface association.

\textbf{Shallow attribution (CWE-327, CWE-798).}
For the MD5 password hash case, the tokeniser splits \texttt{md5} into two subword
tokens \texttt{md} and \texttt{5}; both appear at the top of the attribution ranking
(scores 0.0086 and approximately 0.008 respectively), with all other tokens an order of
magnitude lower.
The prediction is effectively determined by a single algorithm-name pattern.
For hardcoded credentials (CWE-798), \texttt{123} (score 0.0141) and the variable name
\texttt{password} (0.005) dominate: the model identifies the numeric literal assigned to
a security-relevant variable, which is the correct surface signal even if it does not
constitute deep semantic understanding.

Both classes are correctly classified, but for a reason that is fragile: any reformulation
that preserves the vulnerability while removing the identifying token or literal would
likely fool the model.
The ablated CWE-327 case (where \texttt{md5} is replaced with \texttt{hash\_fn})
provides direct evidence of the fallback mechanism: the top tokens shift to
\texttt{password} (0.009), \texttt{encode} (0.008), and \texttt{hexdigest} (0.001),
showing the model can recover a more structural reading of the pattern when the primary
keyword is absent.
This explains why the token ablation test in Section~\ref{sec:understanding} passed:
the model carries a secondary structural representation alongside the primary keyword signal.

\textbf{Incorrect attribution (CWE-78).}
The command injection false positive is the most analytically important result.
On the vulnerable version (user-controlled input reaching \texttt{system()}), attribution
scores are uniformly near zero across all tokens --- the highest is \texttt{ls} at 0.0024,
and \texttt{system} does not appear in the top eight.
The model is confident (p\,=\,0.959) but cannot be said to have identified the
vulnerability; the prediction is too diffuse to attribute to any specific token or
relationship.

On the structurally identical but safe version (hardcoded path, no user input),
\texttt{sprintf} scores 0.671 --- an order of magnitude above every other score
in the entire attribution experiment --- and the model again predicts VULNERABLE (p\,=\,0.929).
This directly explains the persistent minimal-fix failure from Section~\ref{sec:understanding}:
the model learned \texttt{sprintf} as a near-sufficient predictor of the CWE-78 class
during training, because in the synthetic training distribution \texttt{sprintf} appeared
in virtually every vulnerable command injection example.
The model therefore cannot be taught by counterexample to distinguish user-controlled from
hardcoded \texttt{sprintf} arguments without an architectural change that can represent
where a value came from.

\subsection{Relationship to Training Data Co-occurrence}

The attribution pattern across all five CWEs reflects a consistent underlying cause:
the model learns whichever feature in the training data most reliably co-occurs with
the vulnerability label.
When that feature is the true cause of the vulnerability --- as with \texttt{free} in
CWE-416 --- attribution is semantically correct.
When a correlated surface token is a more consistent predictor across training examples
--- as with \texttt{md5} for CWE-327 or \texttt{sprintf} for CWE-78 --- the model
learns the surface token instead, and the attribution exposes that substitution.

This has a direct implication for training data design: hard negatives are necessary
precisely because they break surface co-occurrences that are not causal.
The success of the CWE-327 sha256 fix (now correctly classified as SAFE after hard
negative augmentation) confirms this: introducing safe examples that share the
\texttt{hashlib} token structure but lack the dangerous algorithm name forced the model
to abandon the library-level co-occurrence in favour of the algorithm-level distinction.
The CWE-78 false positive cannot be resolved by the same approach because the
distinguishing feature --- data flow from parameter to buffer --- is not representable
as a token or token combination visible to a sequence encoder.
